Monte Carlo Tree Search (MCTS) is a sample-based planning framework for sequential decision making problems, which has produced state-of-the-art results in a wide range of domains. Most of the algorithms underpinning these results are designed for a single objective. Many real-world problems instead involve trade-offs between several conflicting objectives, and the user's preferences over those objectives may not be known at the time of planning.

The maximum entropy objective, a popular exploration bonus, is shown to be misaligned with the standard objective, both in theory and empirically. Boltzmann Tree Search (BTS) is adapted from Maximum ENtropy Tree Search (MENTS) by removing entropy from the objective. Decaying ENtropy Tree Search (DENTS) then re-introduces entropy as a secondary objective in a way that is aligned with the standard objective. BTS and DENTS are shown to be consistent under mild conditions, and outperform Upper Confidence Bound and MENTS on a range of gridworld environments and on the game of Go.

In the multi-objective setting, two Multi-Objective MCTS (MOMCTS) frameworks are introduced for the decision support scenario with linear utilities. The first is Convex Hull MCTS (CHMCTS), which generalises prior MOMCTS work to stochastic environments via convex hull backups, but these are computationally expensive. The second is the simplex map, an adaptive data structure that maintains value and entropy estimates as a function of the weight vector over the unit simplex via a hierarchical triangulation. Simplex maps integrate cleanly with the entropy bonus used by DENTS, lifting entropy-based exploration to the multi-objective setting. The methods developed in this thesis outperform existing MOMCTS methods on multi-objective gridworld environments and on environments from the MO-Gymnasium library, with simplex map algorithms scaling significantly better than CHMCTS in stochastic environments.